% SOURCE_AUTHORITY: sga5_fr_workpass.tex, lines 6328--6385 (LF-normalized unit).
% SOURCE_SHA256: 791F4EFFC5E02832D5D77ED03518C8156D6F07E4C8238B03545DB93D883FBB28
\subsection*{6.18}
Pongamos
\begin{equation}
N_G=\sum_{g\in G} g\in\Lambda[G].
\tag{6.18.1}
\end{equation}
Sea $X$ un $S$-esquema. Para $L\in\Ob D({}_{\Lambda[G]}X)$, la multiplicación a la izquierda por $N_G$ induce una flecha de $D(X)$
\begin{equation}
N_G:\Lambda\lten_{\Lambda[G]}L
\longrightarrow
\uRHom_{\Lambda[G]}(\Lambda,L),
\tag{6.18.2}
\end{equation}
donde $\Lambda$ se considera como una $\Lambda[G]$-álgebra mediante la aumentación natural, que envía todo $g\in G$ a $1$.

\begin{statement}{Lema 6.18.3}
Si $L\in\Ob D_{\mathrm{ctf}}({}_{\Lambda[G]}X)$, la flecha 6.18.2 es un isomorfismo.
\end{statement}

En efecto, por devissage nos reducimos a suponer primero que $L$ es perfecto, y después que $L=\Lambda[G]$; la afirmación se deduce entonces de que $N_G$ (6.18.1) es una base del módulo de elementos invariantes bajo la acción de $G$ del $\Lambda[G]$-módulo izquierdo $\Lambda[G]$.

\begin{statement}{Proposición 6.19}
En la situación de 6.11, sean $M\in\Ob D_{\mathrm{ctf}}(Y)$,
\[
v\in\Hom(d_1^*M,d_2^!M),
\]
y $a:M\to f_*L$ una flecha de $D({}_{\Lambda[G]}Y)$, considerándose $M$ como objeto de $D({}_{\Lambda[G]}Y)$ mediante la aumentación $\Lambda[G]\to\Lambda$, tales que el cuadrado
\[
\begin{tikzcd}[column sep=large,row sep=large]
d_1^*M \arrow[r,"v"] \arrow[d,"d_1^*a"'] & d_2^!M \arrow[d,"d_2^!a"] & {}\\
d_1^*f_*L \arrow[r,"f_*u"'] & d_2^!f_*L
\end{tikzcd}
\]
sea conmutativo. Se supone que $f_*M\in\Ob D_{\mathrm{ctf}}({}_{\Lambda[G]}Y)$ y que $a$ induce un isomorfismo
\[
M\xrightarrow{\sim}\uRHom_{\Lambda[G]}(\Lambda,f_*L).
\]
Entonces se tiene
\begin{equation}
\Tr_{\Lambda}(v)=
\sum_{g\in G_{\natural}}\Tr_{\Lambda[G]}(f_*u)(g^{-1}).
\tag{6.19.1}
\end{equation}
\end{statement}

Según 6.18.3, la flecha $a$, seguida de $N_G^{-1}$, proporciona un isomorfismo
\[
M\xrightarrow{\sim}\Lambda\lten_{\Lambda[G]}f_*L,
\]
mediante el cual $v$ se identifica con
\[
\Lambda\lten_{\Lambda[G]}f_*u.
\]
Según 6.7.2, $\Tr_{\Lambda}(v)$ es, por tanto, la imagen de $\Tr_{\Lambda[G]}(f_*u)$ mediante la flecha inducida por la aumentación
\[
\Lambda[G]\longrightarrow \Lambda;
\]
ahora bien, esta, por definición, envía cada clase de conjugación a $1$, de donde resulta la conclusión.
